\chapter{Heat Balance simulation}

The theoretical model of the heat balance problem can be found in Ref.\cite{wachutka}. 
\section{Heat equation}
The steady state heat equation with the continuity equation reads as

\begin{equation}\label{heat_law}
\nabla \cdot\left(- \kappa \nabla T + \mathbf{J_{S}} \right) = H_{S}
\end{equation}

where $T$ is the temperature and $\kappa$ is the thermal conductivity tensor.
The source terms are rapresented by $\mathbf{J_S}$ and $H_S$ which are, respectively, the heat flux source and the heat source due to other physical subsystems.

\section{Physical model}

Negleting particle effects, the thermal conductivity is only due to the lattice contribution. The lattice thermal conductivity is read from the database.

The thermal model is tagged as \textit{thermal}. In \textbf{options} subsection we indicate the simulation name (\textit{simulation\_name = whatever\_you\_want}) and the simulation domain (\textit{physical\_regions = wherever\_you\_want})


\begin{verbatim}

model thermal
 {

  options
  { 
    simulation_name = whatever_you_want
    physical_regions = wherever_you_want
  }

    ...
}

\end{verbatim}


\subsection{Electron and hole dissipations}


Electron and hole dissipations give the following sources

\begin{equation}\label{Jp}
\mathbf{J_{S}} = - T P_n \mathbf{J_n} -  T P_p \mathbf{J_p}
\end{equation}

\begin{equation}\label{Hp}
H_{S} = - \mathbf{J_n} \cdot \nabla \varphi_n - \mathbf{J_p} \cdot \nabla \varphi_p
\end{equation}

where $P_n$ and $P_p$ are the electron and hole thermoelectric power, respectively.


The physical model \textbf{heat\_source} includes a specific source identified by its \textbf{model} name.
Concerning electron and hole dissipations the model is

 
named \textit{drift\_diffusion\_dissipation}, as reported below.

\begin{verbatim}

   physical_model heat_source
   {		
     model = drift_diffusion_dissipation

     drift_diffusion_simulation = dd_simul_name
   }

\end{verbatim}

In order to include such a heat source we have to use a drift diffusion simulation.

The syntax \textit{drift\_diffusion\_simulation = dd\_simul\_name} allows this connection.

\subsection{Boundary conditions}


The boundary conditions can fix the temperature at a specific contact. This Dirichlet like condition can be interpreted as a heat reservoir and is identified with \textit{type = heat\_reservoir}. Below an example is reported.

\begin{verbatim}





BC_Regions
{
  BC_Region name_BC_region
  {
    type = heat_reservoir
    BC_reg_numb = 3
    temperature = 300
  }
}
\end{verbatim}



With the notation \textit{temperature = 300} we fix the temperature (300 K in this case) at the contact associated to \textit{BC\_reg\_numb = 3}.


\section{Output data}

In this section we describe the \textit{thermal} quantities that can be plotted.
The plot keywords are \textit{T} (node-based data) and \textit{HeatSource} (element-based data), i.e.

\begin{verbatim}
plot = (......,T,HeatSource,.....)
\end{verbatim}
  
The keyword \textit{T} is referred to the temperature data (node-based data).
The keyword \textit{HeatSource} allows to plot all heat sources considerated in the simulation. In the table we report all sources associated to drift diffusion simulations.



\begin{table}[htb]
\center
\begin{tabular}{l||l||l}
\textit{Expression} & \textit{Heat source} & \textit{Label}\\
\hline\hline
$\frac{\left| \mathbf{J_n} \right| ^2}{\sigma_n}$ & \textit{Electrons Joule's effect} & EJoule\\
$\frac{|\mathbf{J_p}|^2}{\sigma_p}$ & \textit{Holes Joule's effect} & HJoule\\
$qR_{SRH}\left(\phi_p - \phi_n + T \left( P_p - P_n \right) \right)$ & \textit{Recombination effect} &  RecSRH\\
$ -T \mathbf{J_n} \cdot \nabla P_n $ & \textit{Electrons Peltier-Thomson effect} & EPelTh \\
$ -T \mathbf{J_p} \cdot \nabla P_p $ & \textit{Holes Peltier-Thomson effect} & HPelTh
\end{tabular}
\caption{Drift diffusion heat sources}
\label{table:srh_params_db}
\end{table}
